% The Focus Economy: Bounded Autonomy Through Incentive-Aligned Resource Management
% Michael Clark — Are-Self Research
\documentclass[man,12pt,floatsintext]{apa7}

\usepackage[utf8]{inputenc}
\usepackage[T1]{fontenc}
\usepackage{amsmath,amssymb,amsfonts}
\usepackage{graphicx}
\usepackage{booktabs}
\usepackage{hyperref}
\usepackage{algorithm}
\usepackage{algpseudocode}
\usepackage{listings}
\usepackage{xcolor}

\definecolor{codebg}{rgb}{0.95,0.95,0.97}
\definecolor{codegreen}{rgb}{0.2,0.6,0.2}
\definecolor{codegray}{rgb}{0.5,0.5,0.5}
\definecolor{codepurple}{rgb}{0.58,0,0.82}

\lstset{
  backgroundcolor=\color{codebg},
  basicstyle=\ttfamily\small,
  breaklines=true,
  captionpos=b,
  commentstyle=\color{codegreen},
  keywordstyle=\color{codepurple},
  numberstyle=\tiny\color{codegray},
  stringstyle=\color{codegreen},
  frame=single,
  numbers=left,
}

\usepackage[style=apa,backend=biber]{biblatex}
\addbibresource{../../templates/references.bib}

\title{The Focus Economy: Bounded Autonomy Through Incentive-Aligned Resource Management}
\shorttitle{THE FOCUS ECONOMY}
\author{Michael Clark}
\affiliation{Independent Researcher \\ scipraxian}

\abstract{
Autonomous AI agents face a fundamental tension: they must reason long enough to solve
complex problems, but not so long that they waste resources or loop indefinitely. Existing
approaches use hard turn limits or token budgets, which are blunt instruments that ignore
the quality of the agent's work. We introduce the Focus Economy, a resource management
system that aligns an agent's incentive structure with productive behavior. Each reasoning
session begins with a focus budget that decreases with every turn and tool call. Crucially,
focus is \emph{restored} when the agent forms novel memories---saving facts with low
cosine similarity to existing knowledge. This creates a self-regulating system where
agents that learn are rewarded with more reasoning capacity, while agents that loop or
produce redundant output are naturally constrained. We formalize the Focus Economy's
dynamics, analyze its equilibrium properties, and demonstrate its effectiveness in
preventing degenerate behavior across diverse reasoning tasks on consumer hardware.
}

\keywords{bounded autonomy, resource management, incentive alignment, autonomous agents,
focus economy, reinforcement via memory}

\begin{document}
\maketitle

\section{Introduction}

The promise of autonomous AI agents is that they can reason, plan, and act without
continuous human supervision. The peril is that unsupervised agents can consume unbounded
resources: token budgets, API costs, compute time, and human attention when they produce
voluminous but unhelpful output.

Current approaches to bounding agent execution fall into two categories: hard limits
(maximum turns, token ceilings) and soft heuristics (confidence thresholds, self-reflection
loops). Hard limits are blunt---they terminate productive reasoning as readily as
unproductive loops. Soft heuristics are unreliable---the same LLM that enters a loop
may also misjudge when to stop.

We propose a third approach: the Focus Economy, where the agent's execution budget is
dynamically linked to the \emph{novelty} of its outputs. An agent that learns new things
earns the right to keep reasoning. An agent that recycles known information depletes its
budget and must conclude.

\subsection{Design Intuition}

The Focus Economy mirrors biological attention: a person engaged in novel learning
maintains focus naturally, while repetitive tasks induce fatigue. We operationalize this
intuition through three mechanisms:

\begin{enumerate}
  \item \textbf{Focus Budget}: A per-session numeric resource that decrements each turn
  \item \textbf{Memory Novelty Rewards}: Focus restoration proportional to the novelty
        of saved memories (measured by cosine distance from existing knowledge)
  \item \textbf{Efficiency Bonus}: Additional focus and XP for concise output
\end{enumerate}

\section{Related Work}

\subsection{Token and Turn Budgets}

Most agent frameworks impose fixed limits. LangChain's \texttt{max\_iterations}
parameter, AutoGen's \texttt{max\_turns}, and similar mechanisms prevent runaway
execution but cannot distinguish productive from unproductive turns.

\subsection{Self-Reflection and Meta-Reasoning}

Reflexion \parencite{} and similar approaches have the agent evaluate its own progress.
This relies on the same model that may be stuck to recognize that it is stuck---a
bootstrapping problem the Focus Economy avoids by using an external signal (memory
novelty) rather than self-assessment.

\subsection{Reward Shaping in RL}

The Focus Economy's incentive structure parallels reward shaping in reinforcement learning,
where auxiliary rewards guide exploration. Unlike RL reward shaping, our mechanism operates
at the application level without requiring model fine-tuning.

\section{Formal Model}

\subsection{Focus Budget Dynamics}

Let $F_0$ denote the initial focus budget for a session:

\begin{equation}
  F_0 = 10 + 0.5 \times (L - 1)
\end{equation}

where $L = \lfloor \text{XP} / 100 \rfloor + 1$ is the identity's current level.

At each turn $t$, focus is consumed:

\begin{equation}
  F_{t+1} = F_t - c_{\text{turn}} - \sum_{i} c_{\text{tool}_i} + r_t
\end{equation}

where $c_{\text{turn}}$ is the per-turn cost, $c_{\text{tool}_i}$ is the cost for each
tool call (drawn from the tool's \texttt{focus\_modifier}), and $r_t$ is the novelty
reward.

\subsection{Novelty Reward}

When the agent saves a memory (Engram), the Hippocampus computes cosine similarity
$s$ against existing engrams. The novelty reward is:

\begin{equation}
  r = \begin{cases}
    \alpha \times (1 - s) & \text{if } s < 0.90 \\
    0 & \text{if } s \geq 0.90 \text{ (duplicate rejected)}
  \end{cases}
\end{equation}

where $\alpha$ is a scaling constant. Memories that are highly novel ($s \ll 0.90$)
restore more focus than those that narrowly clear the deduplication threshold.

\subsection{Efficiency Bonus}

If the previous turn's output length $\leq L \times 1000$ characters:

\begin{equation}
  F_{t+1} \mathrel{+}= 1, \quad \text{XP} \mathrel{+}= 5
\end{equation}

This encourages concise reasoning---agents that say less but learn more are rewarded
twice.

\subsection{Termination}

When $F_t \leq 0$, the session must conclude. The agent is given one final turn to
call \texttt{mcp\_done()} with its conclusion.

\section{Equilibrium Analysis}

% TODO: Analyze steady-state behavior
% - Under what conditions does an agent sustain indefinite reasoning?
% - What task types naturally exhaust focus?
% - How does the novelty threshold interact with knowledge saturation?

\subsection{Sustainability Condition}

A session is self-sustaining when the expected novelty reward per turn exceeds the
expected focus cost:

\begin{equation}
  \mathbb{E}[r_t] > c_{\text{turn}} + \mathbb{E}\left[\sum_i c_{\text{tool}_i}\right]
\end{equation}

In practice, this condition holds early in a session (many novel facts to discover) and
fails as the agent's knowledge saturates---a natural stopping criterion.

\subsection{Degenerate Behavior Prevention}

We identify three degenerate patterns and show how the Focus Economy addresses each:

\begin{enumerate}
  \item \textbf{Infinite loops}: Focus depletes monotonically when no novel memories are
        formed, guaranteeing termination.
  \item \textbf{Verbose padding}: The efficiency bonus penalizes long output, removing
        the incentive to pad responses.
  \item \textbf{Memory spam}: Cosine similarity deduplication rejects near-duplicate
        engrams, preventing agents from gaming the novelty reward.
\end{enumerate}

\section{Implementation}

\subsection{Integration with the Reasoning Loop}

The Focus Economy is embedded in the Frontal Lobe's reasoning session loop. Focus is
tracked as a field on the \texttt{ReasoningSession} model. The Parietal Lobe's tool
execution pathway reports focus modifications back to the session. The Hippocampus's
engram save path computes novelty and returns the reward.

\subsection{XP and Leveling}

Experience points accumulate across sessions on the \texttt{IdentityDisc}. Every 100 XP
grants a level, which increases the starting focus budget. This creates a longitudinal
incentive: identities that consistently learn earn greater autonomy over time.

\section{Evaluation}

% TODO: Empirical evaluation
% - Compare Focus Economy vs. fixed turn limits across task types
% - Measure token efficiency (useful output per token spent)
% - Analyze memory novelty distributions across sessions
% - Case studies: coding tasks, research tasks, creative tasks

\section{Discussion}

\subsection{Implications for AI Safety}

The Focus Economy provides a \emph{mechanical} bound on agent autonomy that does not
depend on the agent's own judgment. This is significant for safety: the bound holds
regardless of prompt injection, jailbreaking, or model capabilities.

\subsection{Limitations}

The system assumes that memory novelty correlates with task progress. In domains where
the agent must reason extensively without producing storable facts (e.g., mathematical
proof search), the Focus Economy may prematurely terminate productive reasoning.

\subsection{Future Work}

Domain-specific focus modifiers that adjust costs based on task type. Multi-agent focus
pooling where agents in a swarm share a collective budget. Adaptive novelty thresholds
that relax as knowledge domains expand.

\section{Conclusion}

The Focus Economy demonstrates that bounded autonomy need not be blunt. By linking
execution budgets to the novelty of an agent's learning, we create a self-regulating
system that naturally distinguishes productive reasoning from wasteful loops. The
mechanism requires no model fine-tuning, operates at the application layer, and
integrates naturally with existing LLM agent architectures.

\printbibliography

\end{document}
